\DocumentMetadata{
  pdfversion=2.0,
  pdfstandard=ua-2,
  lang=en-US,
  tagging=on,
  tagging-setup={math/setup=mathml-SE,tabsorder=structure}
}
\documentclass[11pt,letterpaper]{article}
\usepackage[margin=0.68in,includefoot]{geometry}
\usepackage{newcomputermodern}
\usepackage{amsmath,amscd,mathtools}
\usepackage{tikz,adjustbox}
\providecommand{\square}{\mdlgwhtsquare}
\usepackage[hidelinks]{hyperref}
\hypersetup{
  pdftitle={Algebra PhD Exam, September 1993},
  pdfauthor={University of Florida Department of Mathematics},
  pdfsubject={Graduate mathematics examination},
  pdfdisplaydoctitle=true,
  bookmarksopen=true
}
\setcounter{secnumdepth}{0}
\setlength{\parindent}{0pt}
\setlength{\parskip}{0.28em}
\setlength{\emergencystretch}{3em}
\setlength{\leftmargini}{2.15em}
\setlength{\leftmarginii}{2.65em}
\setlength{\leftmarginiii}{3em}
\renewcommand{\labelenumi}{\textbf{\theenumi.}}
\pagestyle{empty}
\raggedbottom
\allowdisplaybreaks
\newenvironment{examproblems}[1]{%
  \begin{enumerate}%
  \setcounter{enumi}{#1}%
  \setlength{\topsep}{0.35em}%
  \setlength{\partopsep}{0pt}%
  \setlength{\itemsep}{0.48em}%
  \setlength{\parsep}{0.18em}%
}{\end{enumerate}}
\newenvironment{examsubparts}{%
  \begin{enumerate}%
  \setlength{\topsep}{0.18em}%
  \setlength{\partopsep}{0pt}%
  \setlength{\itemsep}{0.16em}%
  \setlength{\parsep}{0.1em}%
}{\end{enumerate}}
\newenvironment{examinstructions}{%
  \par\begingroup\itshape
}{%
  \par\endgroup\vspace{0.18em}
}
\makeatletter
\renewcommand\section{\@startsection{section}{1}{0pt}%
  {-1.25ex plus -.25ex minus -.1ex}%
  {0.55ex plus .1ex}%
  {\normalfont\large\bfseries}}
\renewcommand{\@maketitle}{%
  \newpage\null\vskip -1em%
  {\footnotesize\bfseries
    \href{https://www.ufl.edu/}{University of Florida}, \href{https://math.ufl.edu/}{Department of Mathematics}\par}%
  \vskip -0.5em%
  \tagstructbegin{tag=Title}%
    {\LARGE\bfseries\href{https://gma.math.ufl.edu/exams/algebra-phd/}{Algebra PhD Exam}, September 1993\par}%
  \tagstructend%
  \vskip 0.25em%
  \tagmcbegin{artifact}%
    \hrule height 1pt%
  \tagmcend%
  \vskip 1em%
}
\makeatother
\title{Algebra PhD Exam, September 1993}
\author{}
\date{}
\begin{document}
\maketitle
\thispagestyle{empty}
\begin{examinstructions}
Do seven of the following questions (the exam committee will NOT select your best seven answers if you answer more than seven). You may quote standard results (within reason) as long as you make it clear that you are doing so and you state them clearly.
\end{examinstructions}
\begin{examproblems}{0}
\item
This problem has two parts.
\begin{examsubparts}
\item[\textnormal{(i)}]
Define what is meant by a group defined by generators and relations.
\item[\textnormal{(ii)}]
Let~\(D\) be the group defined by generators and relations as follows. It is generated by two elements~\(t\) and~\(s\) that satisfy \(t^2=s^2=1\). Prove that~\(D\) is infinite but every proper homomorphic image of~\(D\) is finite.
\end{examsubparts}
\item
State and prove P. Hall’s Theorem that generalizes Sylow’s for finite solvable groups.
\item
Let~\(q\) be any positive integer. Prove that there exists a field with~\(q\) elements if and only if~\(q\) is a power of a prime. Show furthermore that if~\(q\) is a power of a prime then there is, up to isomorphism, only one field with~\(q\) elements.
\item
Let~\(p\) be a prime number, and \(G=T(n,p)\) be the group of all invertible \(n\times n\) matrices which are lower triangular, over the field \(\mathbb{F}_p\) with~\(p\) elements.
\begin{examsubparts}
\item[\textnormal{(a)}]
Let \(U=\{a\in G:a_{ii}=1,\ \text{for all }i=1,\ldots,n\}\). Prove that~\(U\) is nilpotent.
\item[\textnormal{(b)}]
Show that~\(G\) is solvable but, if \(n>1\), not nilpotent.
\end{examsubparts}
\item
Prove: if a Dedekind domain has only a finite number of nonzero prime ideals then it is a principal ideal domain. (Hint: prove first that each prime is principal, use the Chinese Remainder Theorem).
\item
State and prove Hilbert’s Nullstellensatz. (You may use without proof the following Lemma: If~\(F\) is an algebraically closed field extension of a field~\(K\) and~\(I\) is a proper ideal of \(K[x_1,\ldots,x_n]\), then the affine variety \(V(I)\) defined by~\(I\) in \(F^n\) is nonempty.)
\item
Let~\(R\) be a commutative ring with 1. Prove that the set of all nilpotent elements is the intersection of all the prime ideals of~\(R\).
\item
Let~\(R\) be a semisimple ring with 1. Prove that if~\(R\) is Artinian it is also Noetherian. Give a counterexample to show that the converse is false.
\item
Give an example of a ring~\(R\) with 1 and a free~\(R\)-module~\(F\) possessing a submodule which is not free.
\item
Let \(p(x)\) be a non-zero polynomial of degree \(n>0\) in \(\mathbb{Q}[x]\), and let \(r_1,\ldots,r_n\) be its roots. Assume that all the roots are distinct. Consider the Galois group~\(G\) of \(p(x)\) to be a permutation group on the roots, i.e. a subgroup of the symmetric group \(S_n\). Let 
\[
\Delta=\prod_{i<j}(r_i-r_j).
\]
 Prove that \(\Delta^2\in\mathbb{Q}\). Further prove that~\(G\) is contained in the alternating group \(A_n\) if and only if \(\Delta^2\) is a square in~\(\mathbb{Q}\).
\item
Let~\(R\) be an associative ring.
\begin{examsubparts}
\item[\textnormal{(a)}]
Define projective module over~\(R\).
\item[\textnormal{(b)}]
Prove that if \(P=\bigoplus_{\lambda\in\Lambda}P_\lambda\), as~\(R\) modules, then~\(P\) is projective if and only if each \(P_\lambda\) is projective.
\item[\textnormal{(c)}]
Show that any free~\(R\)-module is projective.
\item[\textnormal{(d)}]
Give an example of a projective module that is not free.
\end{examsubparts}
\end{examproblems}
\end{document}
